\section{Threat Model}

\subsection{Threat Model}
\label{sec:threadmodel}

The setup is as follows:
a (remote)
user wants to run a computation over some sensitive data inside of a (local) secure enclave.
We trust the processor and hardware to be correctly implemented and bug free.
We assume the presence of a hardware Root of Trust (RoT) that consists of a secure Read Only Memory (ROM) which contains the Secure Bootloader code that, we assume, cannot be altered.
We also assume the presence of a Physical Unclonable Function (PUF) \cite{herder2014physical} to generate a chip-specific root key or secure non-volatile memory that stores the key.
The rest of the trusted code base is minimal and includes the Secure Bootloader, the Security Monitor (together \textasciitilde15K lines of code (LOC)) and the enclave code.
This means we neither consider attacks that exploit software bugs in the enclave's code nor protect enclaves from outputting their own secrets.
Nevertheless, most of the software stack is considered compromised by the attacker, who can potentially run arbitrary malicious code.
This includes user-level and other enclave programs collocated on the same machine, but also the OS and/or a potential hypervisor.

We are concerned with an attacker mounting software attacks from code collocated on the same machine as our enclave.
This means we ignore any attacks requiring physical access to the machine such as physical (e.g., noise, electromagnetic field, etc.) side channels. 
We also ignore side channels that can be mounted in software, but exploit the physical layer like power side channels.
We do not consider Rowhammer attacks as part of our threat model.
We consider denial-of service attacks \textit{against} the OS and consider a specific threat model where user-level applications and malicious enclaves are allowed to attempt any kind of software-mounted attacks against the OS. However, we do \textit{not} consider denial-of-service attacks \textit{by} the OS.

We consider the usual threat model for secure enclaves and extend it to include collocated transient execution attacks and microarchitectural side channels.
More precisely, we consider \emph{software attacks where the attacker program needs to be running on the same machine as the victim}. 
This means we do not consider attacks that could still be mounted even if the attacker and the victim program were not collocated on the same machine, such as attacks exploiting the timing of public architectural events explicitly instructed by the programmer.

As a result, our platform aims at protecting enclaves against traditional side channel attacks such as cache attacks, but also (almost) all transient execution attacks as described in \cite{canella2019systematic} (including Spectre \cite{Kocher2018spectre, van2018foreshadow, kiriansky2018speculative, chen2019sgxpectre, koruyeh2018spectre, maisuradze2018ret2spec}  and Meltdown \cite{Lipp2018meltdown, stecklina2018lazyfp, van2018foreshadow} variants) 
However, we do not protect against timing completion attacks or attacks using the timing of IO like NETSpectre \cite{schwarz2019netspectre}, as those could be mounted even by an attacker that is not collocated on the platform. 
We leave the mitigation of this kind of attacks to the enclave programmer.

\subsection{Security Properties}
Given the threat model, our high-level security goals for our platform are three-fold: privacy, integrity and availability.

Privacy can be defined as follows:

\noindent\textbf{P1} \emph{Strong Microarchitectural Isolation:} no adversary should be able to exploit microarchitectural side channels to extract secrets from a security domain.

Integrity can be split into sub-properties:

\noindent\textbf{I1} \emph{Authenticity:} no adversary should be able to cheat a remote user into thinking the correct enclave code is running on the remote machine if it is not.
This property implies the integrity of the secure bootloader and the security monitor.

\noindent\textbf{I2} \emph{Runtime Integrity:}
An adversary should not be able to alter the code or data  of a protected security domain once it has been initialized and sealed. This includes the secure bootloader, the security monitor and enclaves.

Finally, OS availability can be defined as follows.

\noindent\textbf{A1} \emph{OS Availability:}  no adversary running a user-level program or enclave should be able to mount a denial-of-service attack on the OS or to deny it access to hardware and software resources, such as memory regions, cores, or the security monitor.

\textbf{Security Goals}
We enforce \textbf{P1} for the secure bootloader, the security monitor, and for any enclave code.
This means the deployed countermeasures are completely transparent to the programmer.
We also enforce properties \textbf{A1}, \textbf{I1} and \textbf{I2}.

We do not protect the OS against microarchitectural side channels, and we are not concerned with an adversary (especially a malicious OS) denying service to our enclaves.

\textbf{Implementation Caveat}
Because our hardware design is loaded on a cloud FPGA without any integrity guarantee, our threat model is not realistic for our prototype and we cannot provide the same security as we potentially could with a silicon chip.
Specifically, we cannot provide any real RoT for our platform like an unmodifiable ROM and a secure PUF.
We still implement the attestation mechanism for demonstration completeness, but replace the PUF by a secret \emph{insecurely} hardcoded in the Bootloader.